% =============================================================================
% Practico 03 - Autoreferencia y Teorema de la Recursion
% Computabilidad y Complejidad (Teoria de la Computacion II)
% FCEFQyN - Universidad Nacional de Rio Cuarto
% Primer Cuatrimestre de 2026
% -----------------------------------------------------------------------------
% Compilacion sugerida: pdflatex resolucion-practico-03.tex
% =============================================================================

\documentclass[11pt,a4paper]{article}

% --- Codificacion y idioma --------------------------------------------------
\usepackage[utf8]{inputenc}
\usepackage[T1]{fontenc}
\usepackage[spanish,es-tabla]{babel}
\usepackage{lmodern}
\usepackage{microtype}

% --- Matematica -------------------------------------------------------------
\usepackage{amsmath,amssymb,amsthm,mathtools}
\usepackage{stmaryrd}

% --- Layout y tipografia ----------------------------------------------------
\usepackage[margin=2.7cm]{geometry}
\usepackage{enumitem}
\setlist[itemize]{topsep=2pt,itemsep=2pt,parsep=0pt}
\setlist[enumerate]{topsep=2pt,itemsep=2pt,parsep=0pt}
\usepackage{parskip}
\usepackage{xcolor}
\usepackage{booktabs}
\usepackage{array}

% --- Codigo -----------------------------------------------------------------
\usepackage{fancyvrb}
\fvset{
  fontsize=\footnotesize,
  frame=single,
  framerule=0.3pt,
  xleftmargin=6pt,
  xrightmargin=6pt
}
\DefineVerbatimEnvironment{pycode}{Verbatim}{}

% --- Hyperref ---------------------------------------------------------------
\usepackage[hidelinks]{hyperref}
\hypersetup{
  pdftitle={Practico 03 - Autoreferencia},
  pdfauthor={Computabilidad y Complejidad - UNRC 2026},
  pdfsubject={Resolucion del Practico 03},
  pdfkeywords={recursion, autoreferencia, SELF, quine, oraculo, Turing reducibility}
}

% --- Ambientes de teoremas --------------------------------------------------
\theoremstyle{plain}
\newtheorem{teorema}{Teorema}[section]
\newtheorem{proposicion}[teorema]{Proposicion}
\newtheorem{lema}[teorema]{Lema}
\newtheorem{corolario}[teorema]{Corolario}

\theoremstyle{definition}
\newtheorem{definicion}[teorema]{Definicion}
\newtheorem{ejemplo}[teorema]{Ejemplo}

\theoremstyle{remark}
\newtheorem*{observacion}{Observacion}
\newtheorem*{idea}{Idea}
\newtheorem*{intuicion}{Intuicion}

% --- Comandos de notacion ---------------------------------------------------
\newcommand{\MT}{\mathrm{MT}}
\newcommand{\TM}{\mathrm{TM}}
\newcommand{\enc}[1]{\langle #1\rangle}
\newcommand{\encMw}{\langle M, w\rangle}
\newcommand{\lang}{\mathcal{L}}
\newcommand{\Sstar}{\Sigma^{*}}
\newcommand{\redm}{\le_{m}}
\newcommand{\redT}{\le_{T}}
\newcommand{\compl}[1]{\overline{#1}}
\newcommand{\qaccept}{q_{\text{accept}}}
\newcommand{\qreject}{q_{\text{reject}}}
\newcommand{\ATM}{A_{\TM}}
\newcommand{\HALT}{\mathit{HALT}_{\TM}}
\newcommand{\ETM}{E_{\TM}}
\newcommand{\MIN}{\mathit{MIN}_{\TM}}
\newcommand{\SELF}{\mathrm{SELF}}
\newcommand{\Pw}[1]{P_{#1}}
\newcommand{\Powerset}{\mathcal{P}}

% --- Cabecera ---------------------------------------------------------------
\title{\textbf{Practico 03 -- Autoreferencia y Teorema de la Recursi\'on}\\[0.3em]
       \normalsize Computabilidad y Complejidad (Teor\'ia de la Computaci\'on II)\\
       \normalsize FCEFQyN -- Universidad Nacional de R\'io Cuarto\\
       \normalsize Primer Cuatrimestre de 2026}
\author{}
\date{}

\begin{document}
\maketitle
\thispagestyle{empty}

\hrule
\vspace{0.4em}
\begin{center}
\small
Resoluci\'on completa del Pr\'actico 03. Apoyada en el Cap\'itulo 6 del libro
de M.~Sipser, \emph{Introduction to the Theory of Computation}, y en las l\'aminas
te\'oricas de la c\'atedra (\texttt{autoref.pdf}).
\end{center}
\vspace{0.4em}
\hrule

\tableofcontents
\bigskip
\hrule
\bigskip

% =============================================================================
\section*{Introducci\'on}
\addcontentsline{toc}{section}{Introducci\'on}

El presente trabajo resuelve la totalidad de los ejercicios del Pr\'actico 03
de la asignatura, centrado en la \emph{autoreferencia} y el \emph{Teorema de la
Recursi\'on}. Estas herramientas permiten que las M\'aquinas de Turing obtengan
su propia descripci\'on y operen con ella, lo cual da lugar a una nueva
demostraci\'on de la indecibilidad de $\ATM$, a propiedades llamativas como
la no-reconocibilidad de $\MIN$, y a la noci\'on de \emph{introspecci\'on}
en lenguajes de programaci\'on reales. La gu\'ia te\'orica utilizada es:

\begin{itemize}
  \item el Cap\'itulo 6 del libro de Sipser (\texttt{chapter-06.pdf}), que
        cubre el Teorema de la Recursi\'on, sus aplicaciones (m\'aquinas
        m\'inimas, teorema del punto fijo) y las nociones de \emph{oracle
        Turing machine} y de \emph{Turing reducibility};
  \item las l\'aminas de la c\'atedra (\texttt{autoref.pdf}) de Pablo Castro,
        que sintetizan la construcci\'on de la m\'aquina $\SELF$, su versi\'on
        recursiva general, el lenguaje $\MIN$ y la noci\'on de or\'aculo.
\end{itemize}

\paragraph{Convenciones de notaci\'on.}
A lo largo del documento adoptamos las siguientes convenciones, en l\'inea
con Sipser:
\begin{itemize}
  \item $\Sigma$ es un alfabeto y $\Sstar$ el conjunto de cadenas finitas
        sobre $\Sigma$.
  \item $\enc{M}$ denota una codificaci\'on (descripci\'on) de la M\'aquina
        de Turing $M$, y $\encMw$ codifica el par $(M,w)$.
  \item $\lang(M)$ es el lenguaje reconocido por $M$.
  \item $M^{O}$ denota una MT $M$ con un or\'aculo para el lenguaje $O$
        (Definici\'on 6.18 de Sipser).
  \item $A\redT B$ significa que $A$ es decidible relativo a $B$, es decir,
        existe una MT con or\'aculo $M^{B}$ que decide $A$
        (Definici\'on 6.20 de Sipser).
  \item ``Decidible'' = Turing-decidible; ``reconocible'' = Turing-reconocible.
\end{itemize}

\paragraph{Lemas de referencia.}
La herramienta b\'asica de toda esta secci\'on es el siguiente lema, que se usa
repetidamente:

\begin{lema}[Sipser 6.1]
\label{lem:q}
Existe una funci\'on computable $q:\Sstar\to\Sstar$ tal que, para toda
$w\in\Sstar$, $q(w)$ es la descripci\'on $\enc{\Pw{w}}$ de una M\'aquina de
Turing $\Pw{w}$ que en cualquier entrada borra la cinta, escribe $w$ y se
detiene.
\end{lema}

A partir de \'el, Sipser construye la m\'aquina $\SELF$ y demuestra el
Teorema 6.3 (Recursi\'on), que es la versi\'on ``parametrizada'' de
$\SELF$.

\bigskip\hrule\bigskip

% =============================================================================
\section{Ejercicio 1: lectura del Cap\'itulo 6}
\label{sec:ej1}

El primer \'item del pr\'actico consiste en \emph{leer el Cap\'itulo 6 del
libro de Sipser}. Se asume realizado como prerrequisito te\'orico para el
resto de los ejercicios. Las nociones que se utilizan a partir del
Ejercicio~2 incluyen, entre otras:
\begin{itemize}
  \item el Lema~\ref{lem:q} (Sipser 6.1) y la construcci\'on de la m\'aquina
        $\SELF$ que imprime su propia descripci\'on;
  \item el Teorema de la Recursi\'on (Sipser 6.3) y su aplicaci\'on a la
        indecibilidad de $\ATM$ (Teorema 6.5);
  \item la no reconocibilidad de $\MIN$ (Teorema 6.7);
  \item la versi\'on de punto fijo de la recursi\'on (Teorema 6.8);
  \item la noci\'on formal de \emph{or\'aculo} y de \emph{Turing reducibility}
        (Definiciones 6.18 y 6.20);
  \item el resultado de cardinalidad: el conjunto de lenguajes es incontable
        (Corolario 4.18) y, por lo tanto, existen lenguajes no reconocibles
        a\'un por MT con or\'aculo (Problema 6.19).
\end{itemize}
Cada ejercicio que sigue cita expl\'icitamente las definiciones y resultados
del cap\'itulo que invoca.

\bigskip\hrule\bigskip

% =============================================================================
\section{Ejercicio 2: la funci\'on \texttt{Self} en Python}
\label{sec:ej2}

\subsection*{Enunciado}

Escribir en Python la funci\'on $\SELF$ que se reproduce a s\'i misma.

\subsection*{Idea y conexi\'on con la teor\'ia}

La m\'aquina $\SELF$ de Sipser est\'a construida en dos partes $A$ y $B$
(p\'aginas 246--248). Por el Lema~\ref{lem:q} tomamos $A = \Pw{\enc{B}}$, la
m\'aquina que imprime $\enc{B}$ en la cinta. Despu\'es, $B$ lee de la cinta
$\enc{B}$, aplica $q$ para reconstruir $\enc{A}$, y combina ambas mitades para
imprimir $\enc{AB}=\enc{\SELF}$.

En un lenguaje de programaci\'on real esta construcci\'on se traduce en un
\emph{quine}: un programa que produce una copia exacta de su c\'odigo fuente.

\subsection*{Implementaci\'on en Python}

\begin{pycode}
# self.py - Quine clasico estilo Sipser
# Parte A: la sentencia print(...) que imprime la cadena B.
# Parte B: la cadena B que describe el programa completo, con un
# marcador {!r} que se reemplaza por la propia B citada.
B = 'B = {!r}\nprint(B.format(B))\n'
print(B.format(B))
\end{pycode}

\paragraph{Correspondencia con $\SELF$.}
La siguiente tabla muestra c\'omo se corresponden las componentes:

\begin{center}
\begin{tabular}{lll}
\toprule
Sipser ($\SELF$)                  & Python (\texttt{self.py})              \\
\midrule
$B$ (cadena base)                 & literal de la variable \texttt{B}      \\
$A = \Pw{\enc{B}}$                & \texttt{print(B.format(B))}            \\
$q$ (computa $\Pw{w}$)            & especificador \texttt{\{!r\}} de \texttt{format} \\
salida $\enc{AB}$                 & texto impreso por \texttt{print}       \\
\bottomrule
\end{tabular}
\end{center}

\paragraph{Verificaci\'on.}
La correcci\'on del quine se comprueba ejecutando
\[
  \texttt{diff <(python3 self.py) self.py}
\]
que no debe producir salida. Esto es porque $\texttt{repr}(B)$ genera el
literal exacto de \texttt{B} que aparece en el c\'odigo, y la cadena
\texttt{B} describe las dos l\'ineas del programa.

\subsection*{Variante mediante introspecci\'on}

Una versi\'on m\'as cercana al enunciado teorico ``una MT obtiene su propia
descripci\'on'' usa el m\'odulo \texttt{inspect}:

\begin{pycode}
import inspect

def self_program():
    """Imprime su propio codigo fuente."""
    print(inspect.getsource(self_program))

if __name__ == "__main__":
    self_program()
\end{pycode}

Aqu\'i \texttt{inspect.getsource} cumple el rol de la operaci\'on
``obtener $\enc{M}$'' que el Teorema 6.3 garantiza para una MT. La l\'inea
\texttt{self\_program()} es la \emph{ejecuci\'on} de esa descripci\'on.

\subsection*{Conclusi\'on}

La existencia del quine en Python es una evidencia pr\'actica del Teorema de
la Recursi\'on: cualquier lenguaje de programaci\'on Turing-completo permite
escribir un programa que produce una copia exacta de s\'i mismo. La
construcci\'on de Sipser para $\SELF$ se traduce directamente en el patr\'on
$A+B$ usando \texttt{format} y \texttt{repr}.

\bigskip\hrule\bigskip

% =============================================================================
\section{Ejercicio 3: dos MT que se imprimen mutuamente}
\label{sec:ej3}

\subsection*{Enunciado}

Describir dos M\'aquinas de Turing $M$ y $N$ tales que, comenzando en
cualquier input, $M$ imprima $\enc{N}$ y $N$ imprima $\enc{M}$.

\begin{observacion}
El enunciado del pr\'actico (\texttt{practica03.pdf}) tiene una errata menor:
pide ``$M$ imprime $\enc{M}$ y $N$ imprime $\enc{M}$'', lo cual es trivial
(dos copias de $\SELF$ resuelven el caso). Tanto el Problema~6.6 de Sipser
como las l\'aminas de c\'atedra (\texttt{autoref.pdf}) piden imprimirse
\emph{mutuamente}, y esa es la versi\'on interesante que resolvemos.
\end{observacion}

\subsection*{Estrategia}

La t\'ecnica es la misma que la de $\SELF$, agregando un \emph{paso de
transformaci\'on} al final. Por el Lema~\ref{lem:q}, dada cualquier cadena
$w$, sabemos construir computacionalmente $\enc{\Pw{w}}$. En particular,
podemos enviar
\[
  \enc{N} \;\longmapsto\; q(\enc{N}) \;=\; \enc{\Pw{\enc{N}}}.
\]
Si definimos $M = \Pw{\enc{N}}$, entonces $q(\enc{N})$ es \emph{exactamente}
$\enc{M}$. Esto convierte ``imprimir $\enc{M}$'' en ``imprimir $q(\enc{N})$'',
una tarea computable a partir de $\enc{N}$.

\subsection*{Construcci\'on}

\paragraph{Definici\'on de $M$.}
$M$ es la m\'aquina del Lema~\ref{lem:q} aplicada a $\enc{N}$:

\medskip\noindent\textbf{$M = $ ``Con entrada $x$:}
\begin{enumerate}
  \item Ignorar $x$.
  \item Borrar la cinta y escribir la cadena $\enc{N}$ (cableada en la
        descripci\'on de $M$).
  \item Detenerse.''
\end{enumerate}
\medskip

Es decir, $M=\Pw{\enc{N}}$, donde $\enc{N}$ aparece como constante en la
descripci\'on de $M$.

\paragraph{Definici\'on de $N$.}
$N$ se construye al estilo $\SELF$, en partes $A_N$ y $B_N$:

\medskip\noindent\textbf{$N$:}
\begin{enumerate}
  \item \textbf{Parte $A_N$.} Es $\Pw{\enc{B_N}}$: ejecutarla deja $\enc{B_N}$
        en la cinta.
  \item \textbf{Parte $B_N$.} En entrada $\enc{B_N}$:
        \begin{enumerate}
          \item computar $q(\enc{B_N}) = \enc{A_N}$;
          \item combinar $A_N$ y $B_N$ para obtener $\enc{A_N B_N}=\enc{N}$;
          \item computar $q(\enc{N})=\enc{M}$;
          \item borrar la cinta, escribir $\enc{M}$ y detenerse.
        \end{enumerate}
\end{enumerate}
\medskip

\subsection*{Por qu\'e no hay circularidad}

A primera vista parecer\'ia que $M$ y $N$ se definen mutuamente y la
construcci\'on es circular. No lo es:
\begin{itemize}
  \item $N$ \emph{calcula} $\enc{N}$ en tiempo de ejecuci\'on (al estilo
        $\SELF$, lo reconstruye desde la cinta), y reci\'en luego produce
        $\enc{M}$ aplicando $q$.
  \item $M$, en cambio, tiene $\enc{N}$ \emph{cableado} como constante; pero
        $\enc{N}$ ya est\'a fijo una vez construida la m\'aquina $N$.
\end{itemize}

\subsection*{Demostraci\'on alternativa por el Teorema de la Recursi\'on}

\begin{teorema}[Sipser 6.3]
Sea $T$ una MT que computa $t:\Sstar\times\Sstar\to\Sstar$. Entonces existe
una MT $R$ que computa $r:\Sstar\to\Sstar$ con
$r(w)=t(\enc{R},w)$ para todo $w$.
\end{teorema}

Aplicando el Teorema con la transformaci\'on $t(\enc{R},w) = q(\enc{R})$
obtenemos directamente una MT $N$ que en cualquier entrada $w$ produce
$q(\enc{N}) = \enc{M}$. La m\'aquina $M$ se obtiene como $\Pw{\enc{N}}$.

\subsection*{Correctitud}

\begin{proof}
\begin{itemize}
  \item Por construcci\'on de $\Pw{w}$ (Lema~\ref{lem:q}), $M$ ignora su
        entrada y deja $\enc{N}$ en la cinta. Luego $M$ imprime $\enc{N}$.
  \item Por la construcci\'on de $\SELF$ (Lema~\ref{lem:q} aplicado a $B_N$)
        $N$ reconstruye su propia descripci\'on $\enc{N}$ en la cinta.
        El \'unico paso adicional es aplicar $q$ una vez m\'as para obtener
        $q(\enc{N})=\enc{M}$. Luego $N$ imprime $\enc{M}$.
\end{itemize}
\end{proof}

\subsection*{Conclusi\'on}

\begin{center}\fbox{$N$ usa la receta de $\SELF$ y aplica $q$ una vez m\'as
para imprimir $\enc{M}$; $M=\Pw{\enc{N}}$.}\end{center}

\bigskip\hrule\bigskip

% =============================================================================
\section{Ejercicio 4: cardinalidad y or\'aculo $\ATM$}
\label{sec:ej4}

\subsection*{Enunciado}

Usar la propiedad de que el conjunto de lenguajes es incontable para demostrar
que existen lenguajes que no pueden ser reconocidos con una m\'aquina con un
or\'aculo $\ATM$.

\subsection*{Idea}

Argumento cl\'asico de cardinalidad (Cantor):

\begin{itemize}
  \item El conjunto de lenguajes sobre $\Sigma$ es \emph{incontable}
        (Corolario 4.18 de Sipser).
  \item El conjunto de MT con or\'aculo $\ATM$ es \emph{contable}: cada una
        se describe con una cadena finita.
  \item La imagen de un conjunto contable bajo cualquier funci\'on es
        contable; en particular, la familia de lenguajes reconocidos por
        MT con or\'aculo $\ATM$ es contable.
  \item Una funci\'on no puede ser sobreyectiva de un conjunto contable a uno
        incontable.
\end{itemize}

Por lo tanto, existen lenguajes que ninguna MT con or\'aculo $\ATM$ reconoce.

\subsection*{Recordatorios}

\begin{teorema}[Sipser, Corolario 4.18]
\label{teo:incont}
El conjunto $\Powerset(\Sstar)$ de lenguajes sobre $\Sigma$ es incontable.
\end{teorema}

\begin{definicion}[Sipser 6.18]
Un \emph{or\'aculo} para el lenguaje $B$ es un dispositivo externo capaz de
responder, en un \'unico paso, si una cadena $w$ pertenece o no a $B$. Una
\emph{MT con or\'aculo}, notada $M^{B}$, es una MT modificada con la
capacidad de hacer consultas al or\'aculo. La descripci\'on $\enc{M^{B}}$
de una MT con or\'aculo es finita (no incluye al or\'aculo en s\'i; solo
contiene los estados de consulta).
\end{definicion}

\subsection*{Demostraci\'on}

\begin{proposicion}
\label{prop:ej4}
Existe un lenguaje $L\subseteq\Sstar$ que no es reconocido por ninguna MT con
or\'aculo $\ATM$.
\end{proposicion}

\begin{proof}
Sea $\mathcal{M}_{\ATM}$ la familia de MT con or\'aculo $\ATM$. La
asignaci\'on
\[
  M^{\ATM} \;\longmapsto\; \enc{M^{\ATM}}\in\Sstar
\]
es una inyecci\'on, ya que distintas MT producen distintas descripciones.
Como $\Sstar$ es contable, $\mathcal{M}_{\ATM}$ tambi\'en lo es.

Definimos
\[
  \mathcal{R}_{\ATM} \;=\; \{\,\lang(M) : M\in\mathcal{M}_{\ATM}\,\}.
\]
La imagen de un conjunto contable bajo cualquier funci\'on es contable, por
lo que $\mathcal{R}_{\ATM}$ es contable.

Por el Teorema~\ref{teo:incont}, $\Powerset(\Sstar)$ es incontable. Luego
\[
  |\mathcal{R}_{\ATM}| \;\le\; \aleph_0 \;<\; 2^{\aleph_0} \;=\;
  |\Powerset(\Sstar)|,
\]
y por lo tanto $\mathcal{R}_{\ATM}\subsetneq \Powerset(\Sstar)$. Existe
$L\in\Powerset(\Sstar)\setminus\mathcal{R}_{\ATM}$; ese $L$ no es reconocido
por ninguna MT con or\'aculo $\ATM$.
\end{proof}

\subsection*{Interpretaci\'on}

\begin{observacion}
Una MT con or\'aculo $\ATM$ es muy poderosa: decide $\ATM$ por consulta
directa, y por construcci\'on adicional decide $\ETM$, $\HALT$, $\mathit{EQ}_{\TM}$, etc.
(Ejemplo 6.19 y ejercicios de Sipser). Sin embargo, el argumento de
cardinalidad muestra que su poder es ``solamente contable'': se quedan
incontables lenguajes fuera. Este resultado es \emph{robusto} bajo el cambio
de or\'aculo: para cualquier lenguaje fijo $B$ existen lenguajes no
reconocibles por MT con or\'aculo $B$.
\end{observacion}

\subsection*{Conclusi\'on}

\begin{center}\fbox{$\exists L\subseteq\Sstar$ tal que $L\notin\{\lang(M^{\ATM}) : M \text{ es MT}\}$.}\end{center}

\bigskip\hrule\bigskip

% =============================================================================
\section{Ejercicio 5: clase \texttt{Persona} e introspecci\'on}
\label{sec:ej5}

\subsection*{Enunciado}

Escribir una clase \texttt{Persona} en Python con los atributos cl\'asicos. Y
hacer una funci\'on que con introspecci\'on:
\begin{itemize}
  \item imprima la lista de atributos de la clase;
  \item imprima el c\'odigo de alguno de los m\'etodos.
\end{itemize}

\subsection*{Contexto te\'orico}

El Teorema de la Recursi\'on dice que toda MT puede obtener su propia
descripci\'on $\enc{M}$ y operar con ella. En un lenguaje de programaci\'on
moderno, esto se concreta como \emph{introspecci\'on}: posibilidad de
consultar, en tiempo de ejecuci\'on, el c\'odigo y la estructura de los
propios objetos. Python lo expone mediante el m\'odulo est\'andar
\texttt{inspect}.

\subsection*{Clase \texttt{Persona}}

\begin{pycode}
# persona.py
class Persona:
    """Representa una persona con datos personales clasicos."""

    especie = "Homo sapiens"  # atributo de clase

    def __init__(self, nombre: str, edad: int, dni: str, email: str):
        self.nombre = nombre
        self.edad = edad
        self.dni = dni
        self.email = email

    def saludar(self) -> str:
        return f"Hola, soy {self.nombre} y tengo {self.edad} anios."

    def es_mayor_de_edad(self) -> bool:
        return self.edad >= 18

    def __str__(self) -> str:
        return f"Persona(nombre={self.nombre!r}, edad={self.edad}, dni={self.dni!r})"
\end{pycode}

Los atributos \texttt{nombre}, \texttt{edad}, \texttt{dni}, \texttt{email} son
de \emph{instancia} y se crean en \texttt{\_\_init\_\_}; \texttt{especie} es
un atributo de \emph{clase}.

\subsection*{Funci\'on de introspecci\'on}

\begin{pycode}
import inspect

def inspeccionar(cls, metodo: str | None = None) -> None:
    """Lista atributos y muestra el codigo de un metodo de la clase."""
    print(f"--- Inspeccion de la clase {cls.__name__} ---")

    miembros = inspect.getmembers(cls)
    atributos = [n for n, _ in miembros if not n.startswith("__")]
    print("Atributos y metodos publicos:")
    for nombre in atributos:
        print(f"  - {nombre}")

    print("\nClasificacion:")
    for nombre in atributos:
        valor = getattr(cls, nombre)
        tipo = "metodo" if inspect.isfunction(valor) else "atributo de clase"
        print(f"  - {nombre} ({tipo})")

    if metodo is None:
        metodos = [n for n, v in miembros
                   if inspect.isfunction(v) and not n.startswith("__")]
        if not metodos:
            print("\n(No hay metodos publicos para mostrar.)")
            return
        metodo = metodos[0]

    print(f"\nCodigo del metodo `{metodo}`:")
    print(inspect.getsource(getattr(cls, metodo)))


if __name__ == "__main__":
    inspeccionar(Persona)
    inspeccionar(Persona, "es_mayor_de_edad")
\end{pycode}

\paragraph{Explicaci\'on.}
\begin{itemize}
  \item \texttt{inspect.getmembers(cls)} devuelve todos los miembros de la
        clase; filtramos los \texttt{dunder}.
  \item \texttt{inspect.isfunction} permite separar m\'etodos de atributos
        de clase.
  \item \texttt{inspect.getsource} reconstruye, a partir del archivo fuente,
        el texto exacto del m\'etodo. Es el equivalente real al
        ``obtener $\enc{M}$'' del Teorema de la Recursi\'on.
\end{itemize}

\subsection*{Observaciones}

Los atributos de instancia no aparecen al inspeccionar la clase, porque
existen solamente despu\'es de instanciar. Para verlos:
\[
  \texttt{p = Persona("Ada", 36, "0", "ada@x.org"); print(vars(p))}.
\]

\subsection*{Conclusi\'on}

\texttt{inspect} provee de f\'abrica todas las herramientas necesarias para
la introspecci\'on pedida. Su existencia es la traducci\'on concreta, en
Python, de la libertad de una MT para acceder a su propia descripci\'on que
el Teorema de la Recursi\'on garantiza.

\bigskip\hrule\bigskip

% =============================================================================
\section{Ejercicio 6: creaci\'on din\'amica de una clase}
\label{sec:ej6}

\subsection*{Enunciado}

Usando introspecci\'on, hacer una funci\'on que cree una clase din\'amicamente.

\subsection*{Idea}

Python expone \texttt{type(nombre, bases, atributos)} como constructor
universal de clases. Llam\'andolo en tiempo de ejecuci\'on con tres
argumentos crea una clase nueva, igual a la que producir\'ia el azucar
sint\'actico \texttt{class}. Combinado con \texttt{inspect}, podemos
construir clases ``a la carta'' a partir de descripciones, en una analog\'ia
directa con la construcci\'on $q$ del Lema~\ref{lem:q}.

\subsection*{C\'odigo}

\begin{pycode}
# clase_dinamica.py
import inspect
from typing import Callable, Iterable

def crear_clase(
    nombre: str,
    campos: Iterable[str],
    metodos: dict[str, Callable] | None = None,
    bases: tuple[type, ...] = (object,),
):
    """Crea y devuelve una nueva clase usando type(...) y luego la inspecciona."""
    metodos = metodos or {}
    campos = list(campos)

    def __init__(self, *args, **kwargs):
        if len(args) > len(campos):
            raise TypeError(
                f"{nombre}.__init__() recibe a lo sumo {len(campos)} argumentos"
            )
        for c, v in zip(campos, args):
            setattr(self, c, v)
        for c, v in kwargs.items():
            if c not in campos:
                raise TypeError(f"campo desconocido para {nombre}: {c!r}")
            setattr(self, c, v)

    def __repr__(self):
        partes = ", ".join(f"{c}={getattr(self, c, None)!r}" for c in campos)
        return f"{nombre}({partes})"

    atributos = {
        "__init__": __init__,
        "__repr__": __repr__,
        "__doc__": f"Clase generada dinamicamente con campos: {campos}",
        **metodos,
    }
    cls = type(nombre, bases, atributos)

    print(f"[ok] Clase creada: {cls.__name__}")
    print(f"     Bases       : {tuple(b.__name__ for b in cls.__bases__)}")
    metodos_pub = [n for n, _ in inspect.getmembers(cls, inspect.isfunction)
                   if not n.startswith("__")]
    print(f"     Metodos     : {metodos_pub or '(ninguno)'}")
    print(f"     Firma init  : {inspect.signature(cls.__init__)}")
    return cls


if __name__ == "__main__":
    Punto = crear_clase(
        "Punto",
        campos=["x", "y"],
        metodos={
            "norma": lambda self: (self.x ** 2 + self.y ** 2) ** 0.5,
            "trasladar": lambda self, dx, dy: Punto(self.x + dx, self.y + dy),
        },
    )
    p = Punto(3, 4)
    print("Instancia    :", p)
    print("norma        :", p.norma())
    print("trasladar    :", p.trasladar(1, 1))
\end{pycode}

\subsection*{Conexi\'on con la teor\'ia}

\begin{observacion}
La funci\'on \texttt{crear\_clase} es la versi\'on programable de la
construcci\'on del Lema~\ref{lem:q}: as\'i como $q$ toma una cadena $w$ y
produce $\enc{\Pw{w}}$, nuestra funci\'on toma una descripci\'on (nombre +
bases + diccionario de atributos) y produce una clase. \texttt{type} juega el
rol de la transformaci\'on $q$, e \texttt{inspect} permite verificar el
resultado.
\end{observacion}

\subsection*{Conclusi\'on}

\texttt{type} y \texttt{inspect} son suficientes para crear clases en tiempo
de ejecuci\'on y validar su estructura, completando el ciclo de
\emph{metaprogramaci\'on}. Esto cierra la analog\'ia entre el Teorema de la
Recursi\'on (capacidad de operar sobre la propia descripci\'on) y las
facilidades est\'andar de Python para inspecci\'on y construcci\'on
din\'amica de objetos.

\bigskip\hrule\bigskip

% =============================================================================
\section*{Conclusiones}
\addcontentsline{toc}{section}{Conclusiones}

A lo largo del Pr\'actico 03 hemos puesto en pr\'actica el \emph{Teorema de
la Recursi\'on} y sus aplicaciones:

\begin{itemize}
  \item Los Ejercicios~2 y~3 construyen autoreferencia: una MT (y su an\'alogo
        en Python) que se imprime a s\'i misma, y un par de MT que se imprimen
        \emph{mutuamente}. Ambos siguen la receta del Lema~\ref{lem:q}.
  \item El Ejercicio~4 muestra que, aun aument\'andole a una MT el poder de
        un or\'aculo $\ATM$, sigue existiendo una cantidad incontable de
        lenguajes que escapa a su alcance. El argumento es de cardinalidad
        (Cantor) y es robusto bajo cambios de or\'aculo.
  \item Los Ejercicios~5 y~6 traen el Teorema de la Recursi\'on al terreno
        pr\'actico: el m\'odulo \texttt{inspect} y la funci\'on \texttt{type}
        de Python son la versi\'on operativa de ``obtener la propia
        descripci\'on'' y ``construir nuevas m\'aquinas a partir de
        descripciones''.
\end{itemize}

En conjunto, los resultados consolidan la imagen de la autoreferencia como
una herramienta central para entender los l\'imites y las capacidades del
modelo de c\'omputo, y como un puente concreto entre la teor\'ia y la
programaci\'on real.

\end{document}
